% chapters/discussion.tex

\chapter{Discussion}
\label{ch:discussion}

% ── Comparison with BKV 2009 ──────────────────────────────────────
\section{Comparison with Benton--Kennedy--Varming}
\label{sec:disc-bkv}

% TODO: What was in BKV 2009 vs. what we added
% - They had: basic CPO, PCF syntax/semantics (Coq 8.2, no HB)
% - We added: HB hierarchy, enriched categories, nat trans, Yoneda,
%   domain equations, Scott topology, function space CCC, adequacy proof

% ── Design trade-offs ─────────────────────────────────────────────
\section{Design Trade-offs}
\label{sec:disc-tradeoffs}

% TODO: Classical vs. constructive (ClassicalEpsilon in Lift.v, Sums.v)
% TODO: Plain records vs. HB for cont_fun/mono_fun
% TODO: The bilimit_exists axiom

% ── The lift monad obstruction ────────────────────────────────────
\section{The Coinductive Lift Monad}
\label{sec:disc-lift-monad}

% Source: theory/LiftMonad.v
% TODO: Why the delay monad does NOT form a CPO
% The CPO obstruction theorem

% ── Quantum extensions ────────────────────────────────────────────
\section{Toward Quantum Programming Language Semantics}
\label{sec:disc-quantum}

% Source: quantum/*.v context files, KLM2024

\subsection{Quantum CPOs}
\label{sec:disc-qcpo}

% TODO: qCPO definition from Kornell-Lindenhovius-Mislove
% Comparison with classical CPOs

\subsection{Quantum Channels as Morphisms}
\label{sec:disc-quantum-morphisms}

% TODO: CPTP maps, quantum continuity

\subsection{Quantum Enrichment}
\label{sec:disc-quantum-enrichment}

% TODO: qCPO-enriched categories as semantic setting for quantum languages

% ── Limitations ───────────────────────────────────────────────────
\section{Limitations}
\label{sec:disc-limitations}

% TODO:
% - Universe polymorphism not fully audited
% - bilimit_exists is axiomatized
% - Quantum portion is exploratory/future work
% - No extraction to executable code
